%!TEX program = xelatex
%!TEX encoding = UTF-8
\def\ano{2026}
\def\semestre{1}
\def\disciplina{Geometria Analítica}
\def\turma{1}
\def\autor{José Antônio O. Freitas}
\def\instituto{MAT-UnB}

\documentclass{beamer}
\usetheme{Boadilla}
\usecolortheme{beaver}
% \mode<presentation>
\usepackage[singlelinecheck=false]{caption}
\usepackage{lmodern}
\usepackage{textpos}
\usepackage{amssymb}
\usepackage{amsmath,amsfonts,amsthm,amstext}
\usepackage{cases}
\usepackage[justification=centering]{caption}
\usepackage[brazilian]{babel}
% \usepackage[latin1]{inputenc}
\usepackage{graphicx}
\graphicspath{{/home/jfreitas/GitHub/geometria_analitica/apresentacoes/logo}{/home/jfreitas/GitHub/geometria_analitica/notas_de_aula/figuras-latex}{C:/Users/jfreitas/GitHub/geometria_analitica/apresentacoes/logo}}
\usepackage{enumitem}
\usepackage{multicol}
\usepackage{answers}
\usepackage{qrcode}
\usepackage{tikz,ifthen}
\usetikzlibrary{lindenmayersystems}
\usetikzlibrary[shadings]

\newcounter{saveenumi}
\newcommand{\seti}{\setcounter{saveenumi}{\value{enumi}}}
\newcommand{\conti}{\setcounter{enumi}{\value{saveenumi}}}

\resetcounteronoverlays{saveenumi}


\renewcommand{\qedsymbol}{$\blacksquare$}

\setbeamertemplate{itemize items}[ball]

\input{../configuracoes_latex/ambientes}
\newcommand{\nsub}{\varsubsetneq}
\newcommand{\vaz}{\emptyset}
\newcommand{\dom}{{\rm dom\,}}
\newcommand{\im}{{\rm Im\,}}
\newcommand{\sub}{\subseteq}
\newcommand{\n}{\mathbb{N}}
\newcommand{\z}{\mathbb{Z}}
\newcommand{\rac}{\mathbb{Q}}
\newcommand{\real}{\mathbb{R}}
\newcommand{\complex}{\mathbb{C}}
\newcommand{\cp}[1]{\mathbb{#1}}
\newcommand{\ch}{\mbox{\textrm{car\,}}\nobreak}
\newcommand{\vesp}[1]{\vspace{ #1  cm}}
\newcommand{\compcent}[1]{\vcenter{\hbox{$#1\circ$}}}
\newcommand{\comp}{\mathbin{\mathchoice
{\compcent\scriptstyle}{\compcent\scriptstyle}
{\compcent\scriptscriptstyle}{\compcent\scriptscriptstyle}}}
\newcommand\T{\rule{0pt}{2.6ex}}
\newcommand{\elipse}[2]{\dfrac{x^2}{#1} + \dfrac{y^2}{#2} = 1}
\newcommand{\inner}[2]{\langle \vec{#1}, \vec{#2} \rangle}
\newcommand{\proj}{{\rm proj}}
\newcommand{\eq}{=}
\newcommand{\norm}[1]{\left\Vert{#1}\right\Vert}
\newcommand{\integer}{\mathbb{Z}}
\renewcommand{\sin}{{\rm sen\,}}
\renewcommand{\tan}{{\rm tg\,}}
\renewcommand{\csc}{{\rm cossec\,}}
\renewcommand{\cot}{{\rm cotg\,}}
\renewcommand{\sinh}{{\rm senh\,}}




\title{Vetores no Plano}
\author[\autor]{\autor}
\institute[\instituto]{\instituto}
\date{}

\begin{document}
    \begin{frame}
        \maketitle
    \end{frame}

    \logo{
\includegraphics[scale=.1]{logo-MAT.png}\vspace*{8.5cm}}

    \begin{frame}
        Vamos considerar o plano definido pelo par de retas perpendiculares $x$ e $y$ como abaixo:
        \begin{figure}[!h]
            \centering
            \caption{Plano cartesiano}
            \includegraphics{plano-cartesiano.pdf}
        \end{figure}
    \end{frame}

    \begin{frame}
        Seja $P$ um ponto qualquer do plano.\pause 

        Por $P$ podemos traçar uma única reta $r$ paralela à reta $x$.\pause

        E uma única reta $s$ paralela à reta $y$.\pause 

        Estas retas interceptam $x$ e $y$ em pontos $P_x$ \pause e $P_y$, respectivamente.\pause

        Denote por $\alpha$ o número real correspondente a $P_x$ \pause e por $\beta$ o número real correspondente a $P_y$.\pause

        Estes dois números determinam o ponto $P$, \pause isto é, \pause conhecendo $\alpha$ e $\beta$ \pause podemos traças as retas $r$ e $s$.\pause 

        O ponto $P$ é a interseção destas retas.\pause 

        Os números $\alpha$ e $\beta$ \pause são chamados \textbf{abscissa} \pause e \textbf{ordenada} do ponto $P$, \pause respectivamente.\pause

        Eles formam as \textbf{coordenadas} de $P$.
    \end{frame}

    \begin{frame}
        \begin{figure}[!h]
            \centering
            \caption{Coordenadas do ponto $P$}
            \includegraphics{coordenadas-plano-cartesiano.pdf}
        \end{figure}
    \end{frame}

    \begin{frame}

        Denotaremos assim o ponto $P$ por
        \[
            P(\alpha,\beta).
        \]
    \end{frame}

    \begin{frame}


        Sejam $P(x_1,y_1)$ e $Q(x_2,y_2)$ pontos do plano. \pause 

        A partir de $P$ e $Q$ podemos construir o triângulo retângulo $PQS$ \pause como abaixo:
        \begin{figure}[!h]
            \centering
            \caption{Distância entre dois pontos}
            \includegraphics{distancia-pontos-plano-cartesiano.pdf}
        \end{figure}\pause
    \end{frame}

    \begin{frame}
        Este número é chamado \textbf{distância} de $P$ a $Q$ e é indicado por $d(P,Q)$, isto é, por definição
        \[
            d(P,Q) = \sqrt{(x_1 - x_2)^2 + (y_1 - y_2)^2}.\pause
        \]

        Dizemos que dois pontos $P(x_1,y_1)$ e $Q(x_2,y_2)$ \pause são iguais \pause se, e somente se, $x_1 = x_2$ \pause e $y_1 = y_2$.

    \end{frame}

    \begin{frame}
        Sabemos que a cada par ordenado $(\alpha, \beta)$ corresponde um ponto do plano.\pause
        \vspace{.3cm}

        Assim se $(\alpha, \beta) \ne (0,0)$ \pause podemos associar uma seta \pause com início no ponto $O(0,0)$ \pause e final no ponto $(\alpha,\beta)$
        \pause como na figura abaixo:\pause

        \begin{figure}[!h]
            \centering
            \caption{Seta representando o ponto $(\alpha, \beta)$.}
            \includegraphics{vetores-plano-cartesiano.pdf}
        \end{figure}
    \end{frame}

    \begin{frame}
        Assim o par ordenado $(\alpha, \beta)$ pode ser representado graficamente por um ponto ou por uma seta. \pause

        \vspace{.3cm}

        Quando representamos $(\alpha, \beta)$ por uma seta, \pause podemos associar a este para uma \textbf{direção}, \pause dada por uma reta passado por essa
        seta, \pause um \textbf{sentido} \pause, dado pela ponta da seta \pause e um \textbf{módulo},\pause que definiremos a seguir.\\

        \vspace{.3cm}

        O \textbf{módulo} do par $(\alpha, \beta)$ \pause será a distância do ponto $(\alpha, \beta)$ \pause até o ponto $O(0, 0)$ \pause e é dada por\pause
        \[
            \sqrt{(\alpha - 0)^2 + (\beta - 0)^2} = \pause \sqrt{\alpha^2 + \beta^2}\pause
        \]
        que é o comprimento da seta.
    \end{frame}

    \begin{frame}
        Em geral, um objeto ao qual podemos associar os conceitos de direção, \pause sentido \pause e módulo \pause é chamado de \textbf{vetor}.\pause\\

        \vspace{.3cm}
        
        Denotaremos um vetor de extremidade em um ponto $(\alpha, \beta)$ \pause por
        \[
            \vec{u} = (\alpha, \beta).
        \]
    \end{frame}

    \begin{frame}
        Esta notação significa que $\vec{u}$ é a seta com início no ponto $O(0,0)$ \pause e final no ponto $(\alpha, \beta)$.\pause
        \begin{figure}[!h]
            \centering
            \caption{Vetor $ \vec{u} = (\alpha, \beta)$.}
            \includegraphics{vetor-coordenadas-plano-cartesiano.pdf}
        \end{figure}
    \end{frame}

    \begin{frame}
        Por exemplo, $\vec{u} = (3,2)$ \pause e $\vec{v} = (-1,-2)$ \pause são representados graficamente por:\pause
        \begin{figure}[!h]
            \centering
            \caption{Vetores $\vec{u} = (3,2)$ e $\vec{v} = (-1,-2)$.}
            \includegraphics{exemplo-vetores-plano-cartesiano.pdf}
        \end{figure}
    \end{frame}

    \begin{frame}
        O vetor que tem início e fim no ponto $(0,0)$ \pause é chamado de \textbf{vetor nulo} \pause e será denotado por\pause
        \[
            \vec{O} = (0,0).
        \]
    \end{frame}

    \begin{frame}
        Podemos também representar um vetor por uma seta que não parte da origem. \pause Por exemplo, os pontos $A(x_1,y_1)$ e $B(x_2,y_2)$ determinam o
        vetor\pause
        \[
            \vec{AB} = (x_2 - x_1, y_2 - y_1).
        \]
    \end{frame}

    \begin{frame}
        \begin{figure}[!h]
            \centering
            \caption{Vetor $\vec{AB}$}
            \includegraphics{vetor-AB-plano-cartesiano.pdf}
        \end{figure}
    \end{frame}

    \begin{frame}
        Assim a seta que representa o vetor $\vec{AB}$ \pause e a seta que representa o vetor $\vec{u}$, \pause que tem início na origem \pause e extremidade no
        ponto $(x_2 - x_1, y_2 - y_1)$, \pause têm o mesmo módulo, \pause direção \pause e sentido que $\vec{AB}$.\pause\\

        \vspace{.3cm}

        Assim vamos representar o vetor $\vec{AB}$ \pause pela forma que for mais conveniente.
    \end{frame}

    \begin{frame}

        \begin{definicao}
            Seja $\vec{u} = (x_1, y_1)$ um vetor. \pause A \textbf{norma} ou \textbf{módulo} de $\vec{u}$, \pause denotada por $\norm{\vec{u}}$, \pause é dada
            por\pause
          \[
            \norm{\vec{u}} = \pause \sqrt{x_1^2 + y_1^2}.
          \]
        \end{definicao}
    \end{frame}

    \begin{frame}
        Um vetor $\vec{u}$ \pause tal que
        \[
            \norm{\vec{u}} = 1\pause
        \]
        é chamado de \pause \textbf{vetor unitário}.\pause
        \begin{figure}[!h]
            \centering
            \caption{Vetores unitários: $\norm{\vec{u}} = 1$}
            \includegraphics{vetores-unitarios-plano-cartesiano.pdf}
        \end{figure}
    \end{frame}

    \begin{frame}
        \begin{definicao}
            Dois vetores $\vec{u}$ \pause e $\vec{v}$ \pause são \textbf{paralelos} \pause (ou \textbf{colineares}) \pause se as retas que contém estes vetores
            \pause são paralelas (podendo ser coincidentes). \pause Caso $\vec{u}$ e $\vec{v}$ sejam paralelos, \pause escreveremos $\vec{u}\varparallel\vec{v}$.
            \pause Neste caso, $\vec{u}$ e $\vec{v}$ \pause têm a mesma direção, mas podem ter sentidos contrários.
        \end{definicao}
    \end{frame}

    \begin{frame}
        \begin{figure}[!h]
            \centering
            \caption{Vetores paralelos}
            \includegraphics[scale=0.8]{vetores-paralelos-plano-cartesiano.pdf}
        \end{figure}
    \end{frame}

    \begin{frame}
        \begin{definicao}
          Sejam $\vec{u} = (x_1, y_1)$ \pause e $\vec{v} = (x_2, y_2)$ \pause vetores. \pause Então:\pause
          \[
              \vec{u} + \vec{v} = \pause (x_1 + x_2, \pause y_1 + y_2).
          \]
        \end{definicao}
    \end{frame}

    \begin{frame}
        \begin{figure}[!h]
          \centering
          \caption{Soma de vetores}
          \includegraphics{soma-vetores-plano-cartesiano.pdf}
        \end{figure}
    \end{frame}

    \begin{frame}
        \begin{propriedades}
          Sejam $\vec{u}$, \pause $\vec{v}$ \pause e $\vec{w}$ \pause vetores quaisquer. \pause As seguintes propriedades são verdadeiras:\pause
          \begin{enumerate}[label=({\roman*})]
            \item $(\vec{u} + \vec{v}) \pause + \vec{w} \pause = \vec{u} + \pause (\vec{v} + \vec{w})$ \pause (Associatividade da soma)\pause
            \item $\vec{u} + \pause \vec{v} = \pause \vec{v} + \vec{u}$ \pause (Comutatividade da soma)\pause
            \item Existe um único vetor $\vec{w}$ \pause que somado a $\vec{u}$ \pause resulta no próprio vetor $\vec{u}$:\pause
            \[
              \vec{u} + \vec{w} \pause = \vec{w} + \vec{u} \pause = \vec{u}.\pause
            \]
            Tal vetor $\vec{w}$ é o vetor nulo $\vec{0}$. \pause (Elemento Neutro da soma)\pause
            \item Para cada $\vec{u}$, \pause existe um único vetor $\vec{v}$ \pause tal que\pause
            \[
              \vec{u} + \vec{v} \pause = \vec{0} \pause = \vec{v} + \vec{u}.\pause
            \]
          \end{enumerate}
        \end{propriedades}
    \end{frame}

    \begin{frame}
        \begin{definicao}
            O vetor $\vec{v}$ do item \textit{(iv)} \pause da definição anterior \pause é chamado de \textbf{vetor oposto}.
        \end{definicao}
    \end{frame}

    \begin{frame}
        \begin{definicao}
          Sejam $\alpha \in \real$ e \pause $\vec{u} = (x_1, y_1)$ \pause um vetor. \pause Então:\pause
          \[
              \alpha \vec{u} \pause = (\alpha x_1, \pause \alpha y_1).
          \]
        \end{definicao}
    \end{frame}

    \begin{frame}
        \begin{figure}[!h]
          \centering
          \caption{Multiplicação por escalar}
          \includegraphics{multiplicacao-escalar-vetores-plano-cartesiano.pdf}
        \end{figure}
\end{frame}

    \begin{frame}
        \begin{observacao}
            \begin{enumerate}[label=({\roman*})]
                \item Sejam $\vec{u}$ um vetor não-nulo \pause e $\lambda \ne 0$ \pause um escalar. \pause Então $\vec{u} \varparallel \lambda\vec{u}$. \pause
                    Além disso, \pause  se $\lambda > 0$, \pause então $\vec{u}$ e $\lambda\vec{u}$ têm o mesmo sentido, \pause e se $\lambda < 0$, \pause
                    então $\vec{u}$ e $\lambda\vec{u}$ têm sentidos contrários.\pause

                    \vspace{.3cm}

                \item Quando $\lambda = -1$, \pause o vetor $-\vec{u} = (-x_1,-y_1)$ \pause é exatamente o vetor oposto de $\vec{u}$. \pause Note que $-\vec{u}$
                    \pause tem a mesma direção que $\vec{u}$, \pause mas sentido contrário.
            \end{enumerate}
        \end{observacao}
    \end{frame}

    \begin{frame}
        \begin{propriedades}
          Quaisquer que sejam os números reais $\alpha$ \pause e $\beta$ \pause e quaisquer que sejam os vetores $\vec{u}$ \pause e $\vec{v}$ \pause valem as
          igualdades:\pause
          \begin{enumerate}[label=({\roman*})]
            \item $\alpha \pause (\vec{u} + \vec{v}) \pause = \alpha\vec{u} \pause + \alpha\vec{v}$;\pause
            \item $(\alpha + \beta)\pause \vec{u} =\pause \alpha\vec{u} \pause + \beta\vec{u}$;\pause
            \item $1\vec{u} \pause = \vec{u}$;\pause
            \item $\alpha \pause (\beta\pause \vec{u}) \pause = (\alpha\beta)\pause \vec{u} = \beta\pause (\alpha\vec{u})$.
          \end{enumerate}
        \end{propriedades}
    \end{frame}

    \begin{frame}
        \begin{observacao}
            Dado um vetor não nulo $\vec{u}$ o vetor
            \[
                  \vec{v} = \dfrac{\vec{u}}{\norm{\vec{u}}}
            \]
            é um vetor unitário na direção de $\vec{u}$ pois
            \[
                  \norm{\vec{v}} = \norm{\dfrac{\vec{u}}{\norm{\vec{u}}}} = \dfrac{1}{\norm{\vec{u}}}\norm{\vec{u}} = 1.
            \]

        \end{observacao}
    \end{frame}

    \begin{frame}
        \begin{definicao}
            Sejam $\vec{u}$ e $\vec{v}$ vetores não-nulos. Então $\vec{u} = \vec{v}$ se, e somente se, $\vec{u}$ e $\vec{v}$ têm normas iguais, são de mesma direção e de mesmo sentido.
        \end{definicao}

        \begin{proposicao}\label{vetores_paralelos}
          Dois vetores não nulos $\vec{u}$ e $\vec{v}$ são paralelos se, e somente se, existe um escalar $\lambda$ tal que $\vec{u} = \lambda\vec{v}$ (consequentemente, $\lambda \ne 0$ e assim $\vec{v} = \vec{u}/\lambda$).
        \end{proposicao}
    \end{frame}

    \begin{frame}
        \begin{definicao}
          Sejam $\vec{u}$ e $\vec{v}$ vetores não nulos. O \textbf{ângulo} entre $\vec{u}$ e $\vec{v}$ é definido pelo ângulo $\theta$ determinado por $\vec{u}$ e $\vec{v}$ e que satisfaz $0 \le \theta \le \pi$, quando os vetores $\vec{u}$ e $\vec{v}$ são representados com a mesma origem.\index{Vetores!ângulo}
        \end{definicao}
    \end{frame}

    \begin{frame}
        \begin{figure}[!h]
          \centering
          \caption{ângulo entre vetores}
          \includegraphics{angulo-vetores-plano-cartesiano.pdf}
        \end{figure}
    \end{frame}

    \begin{frame}
        \begin{definicao}
          Quando o ângulo $\theta$ entre dois vetores $\vec{u}$ e $\vec{v}$ é reto, isto é, $\theta = \pi/2$, ou um deles é o vetor nulo, dizemos que os vetores $\vec{u}$ e $\vec{v}$ são \textbf{ortogonais} ou \textbf{perpendiculares} entre si. Denotamos tal fato, escrevendo $\vec{u} \perp \vec{v}$.\index{Vetores!Ortogonais}
        \end{definicao}
    \end{frame}

    \begin{frame}
        Sejam $\vec{u}$ e $\vec{v}$ vetores não nulos e $\theta$ o ângulo entre eles.
        \begin{figure}[!h]
          \centering
          \caption{Determinação de $\theta$}
          \includegraphics{determinacao-angulo-vetores-plano-cartesiano.pdf}
        \end{figure}
    \end{frame}

    \begin{frame}
        Aplicando a Lei dos cossenos no triângulo abaixo
        \begin{figure}[!h]
          \centering
          \caption{Lei dos cossenos}
          \includegraphics{lei-cossenos-vetores-plano-cartesiano.pdf}
        \end{figure}
    \end{frame}

    \begin{frame}
        temos
        \begin{align}\label{equacaonorma}
          \norm{\vec{u} - \vec{v}}^2 = \norm{\vec{u}}^2 + \norm{\vec{v}}^2 - 2\norm{\vec{u}}\norm{\vec{v}}\cos\theta.
        \end{align}
    \end{frame}

    \begin{frame}
        \begin{definicao}\label{produtointerno}
          O \textbf{produto escalar} ou \textbf{produto interno} dos vetores $\vec{u}$ e $\vec{v}$, indicado por $\inner{u}{v}$, é o número real tal que:\index{Vetores!Produto Escalar}
          \begin{enumerate}[label=({\roman*})]
            \item Se $\vec{u}$ ou $\vec{v}$ é nulo, então $\inner{u}{v} = 0$.
            \item Se $\vec{u}$ e $\vec{v}$ não são nulos e $\theta$ é o ângulo entre $\vec{u}$ e $\vec{v}$, então
            \begin{align}\label{produto-interno}
              \inner{u}{v} = \norm{\vec{u}}\norm{\vec{v}}\cos\theta.
            \end{align}
          \end{enumerate}
        \end{definicao}
    \end{frame}

    \begin{frame}
        \begin{proposicao}
          Sejam $\vec{u}$ e $\vec{v}$ vetores e $\theta$ o ângulo entre $\vec{u}$ e $\vec{v}$.
          \begin{enumerate}[label=({\roman*})]
            \item Se $\vec{u}$ e $\vec{v}$ não são nulos, então
            \[
              \cos\theta = \dfrac{\inner{u}{v}}{\norm{\vec{u}}\norm{\vec{v}}}.
            \]
            \item Qualquer que seja o vetor $\vec{u}$,
            \[
              \norm{\vec{u}} = \sqrt{\inner{u}{u}}.
            \]
            \item Quaisquer que sejam os vetores $\vec{u}$ e $\vec{v}$, $\vec{u}\perp\vec{v}$ se, e somente se, $\inner{u}{v} = 0$.
            \end{enumerate}
        \end{proposicao}
   \end{frame}

    \begin{frame}
        Agora, como podemos encontrar o produto interno sem necessitar de determinar o ângulo $\theta$ entre dois vetores $\vec{u}$ e $\vec{v}$\,? Substituindo o produto interno na equação \eqref{equacaonorma} temos
        \[
          \norm{\vec{u} - \vec{v}}^2 = \norm{\vec{u}}^2 + \norm{\vec{v}}^2 - 2\inner{u}{v}
        \]
        isto é,
        \begin{align}\label{equacao-auxilar-norma}
          \inner{u}{v} = \dfrac{1}{2}(\norm{\vec{u}}^2 + \norm{\vec{v}}^2 - \norm{\vec{u} - \vec{v}}^2).
        \end{align}
        Se $\vec{u} = (x_1, y_1)$ e $\vec{v} = (x_2, y_2)$, então podemos reescrever \eqref{equacao-auxilar-norma} como
        \begin{align*}
          \inner{u}{v} &= \dfrac{1}{2}(x_1^2 + y_1^2 + x_2^2 + y_2^2 - \norm{(x_1 - x_2, y_1 - y_2)}^2)\\
          &= \dfrac{1}{2}[x_1^2 + y_1^2 + x_2^2 + y_2^2 - (x_1 - x_2)^2 - (y_1 - y_2)^2]\\
          &= x_1x_2 + y_1y_2.
        \end{align*}
    \end{frame}

    \begin{frame}
        Acabamos de demonstrar o seguinte teorema:
        \begin{teorema}
          O produto interno de dois vetores $\vec{u} = (x_1, y_1)$ e $\vec{v} = (x_2, y_2)$ em $\real^2$ é dado por
          \[
            \inner{u}{v} = x_1x_2 + y_1y_2.
          \]
        \end{teorema}
    \end{frame}

    \begin{frame}
        \begin{exemplos}
          \begin{enumerate}[label=({\roman*})]
            \item Sejam $\vec{u} = (1, 1)$ e $\vec{v} = (1,0)$. Determine o ângulo entre $\vec{u}$ e $\vec{v}$.
            \item Encontre um vetor $\vec{u} = (x, y)$ ortogonal a $\vec{v} = (4, -1)$ e tal que $\inner{u}{w} = -1$, onde $\vec{w} = (1,1)$.
          \end{enumerate}
        \end{exemplos}
    \end{frame}

    \begin{frame}
        \begin{proposicao}\label{propriedades-produto-interno}
          Quaisquer que sejam os vetores $\vec{u} = (x_1, y_1)$, $\vec{v} = (x_2, y_2)$ e $\vec{w} = (x_3, y_3)$ e qualquer que seja o número real $\lambda$ temos
          \begin{enumerate}[label=({\roman*})]
            \item\label{linearidade-produto-interno} $\langle \vec{u}, \vec{v} + \vec{w}\rangle = \inner{u}{v} + \inner{u}{w}$
            \item $\langle \vec{u} + \vec{v}, \vec{w}\rangle = \inner{u}{w} + \inner{v}{w}$
            \item\label{linearidade2-produto-interno} $\langle (\lambda \vec{u}\,), \vec{v}\rangle = \langle \vec{u}, (\lambda \vec{v}\,)\rangle = \lambda\inner{u}{v}$
            \item $\inner{u}{v} = \inner{v}{u}$
            \item Se $\vec{u} \ne \vec{0}$, então $\inner{u}{u} > 0$.
          \end{enumerate}
        \end{proposicao}
   \end{frame}

    \begin{frame}
        \begin{observacao}
          \begin{enumerate}[label=({\roman*})]
            \item As propriedades \ref{linearidade-produto-interno} e \ref{linearidade2-produto-interno} Proposição \ref{propriedades-produto-interno} podem ser estendidas para qualquer número de vetores e escalares:
            \[
              \langle\vec{u}, (\lambda_1\vec{v_1} + \lambda_2\vec{v_2} + \cdots + \lambda_n\vec{v_n} )\rangle = \lambda_1\inner{u}{v_1} + \lambda_2\inner{u}{v_2} + \cdots + \lambda_n\inner{u}{v_n}.
            \]
            \item Na igualdade $\inner{u}{v} = \inner{u}{w}$ não podemos concluir que $\vec{v} = \vec{w}$. Por exemplo, para $\vec{u} = (1, 0)$, $\vec{v} = (2, 1)$ e $\vec{w} = (2, -5)$ temos $\inner{u}{v} = \inner{u}{w}$ e no entanto $\vec{v} \ne \vec{w}$. Mas podemos concluir que $\vec{u}\perp(\vec{v} - \vec{w})$.
            \item De $\inner{u}{v} = 0$ não podemos concluir que $\vec{u} = \vec{0}$ ou $\vec{v} = \vec{0}$. Por exemplo, para $\vec{u} = (2, 1)$ e $\vec{v} = (-2, 4)$ temos $\inner{u}{v} = 0$ e no entanto $\vec{u} \ne \vec{0}$ e $\vec{v} \ne \vec{0}$.
          \end{enumerate}
        \end{observacao}
    \end{frame}

    \begin{frame}
        \begin{exemplos}
          Calcule o ângulo $\alpha$ entre $\vec{u} + \vec{v}$ e $\vec{u} - \vec{v}$, sabendo que $\norm{\vec{u}} = \sqrt{5}$, $\norm{\vec{v}} = 1$ e que o ângulo $\theta$ entre $\vec{u}$ e $\vec{v}$ é $\pi/4$.
        \end{exemplos}
    \end{frame}

    \begin{frame}
        \begin{proposicao}\label{DesigualdadeTriangular}
          Sejam $\vec{u}$ e $\vec{v}$ vetores.
          \begin{enumerate}[label=({\roman*})]
            \item $|\inner{u}{v}| \le \norm{\vec{u}}\norm{\vec{v}}$ [Desigualdade de Scharwz]\index{Desigualdade de Scharwz}
            \item $\norm{\vec{u} + \vec{v}} \le \norm{\vec{u}} + \norm{\vec{v}}$ [Desigualdade Triangular]\index{Desigualdade Triangular}
          \end{enumerate}
        \end{proposicao}
    \end{frame}

    \begin{frame}


        Sejam $\vec{u} = (x_1, y_1)$ e $\vec{v} = (x_2, y_2)$ vetores não nulos. Suponha que o ângulo $\theta$ entre $\vec{u}$ e $\vec{v}$ é tal que $0 < \theta < \pi/2$.
        \begin{figure}[!h]
          \centering
          \caption{Projeção ortogonal}
          \includegraphics{projecao-vetores-plano-cartesiano.pdf}
        \end{figure}

        O vetor $\vec{OP}$ é paralelo ao vetor $\vec{u}$ e o vetor $\vec{PA}$ é ortogonal ao vetor $\vec{u}$.
    \end{frame}

    \begin{frame}
        \begin{definicao}\label{projecao_ortogonal}
          Seja $\vec{u}$ um vetor não nulo. Dado qualquer vetor $\vec{v}$, o vetor $\vec{OP}$ é chamado de \textbf{projeção ortogonal} de $\vec{v}$ sobre $\vec{u}$, e é indicado por $\proj_{\vec{u}}\vec{v}$ se satisfaz as seguintes condições:\index{Vetores!Projeção ortogonal}
          \begin{enumerate}[label=({\roman*})]
            \item $\vec{OP} \varparallel\vec{u}$
            \item $(\vec{v} - \vec{OP}) \perp \vec{u}$.
          \end{enumerate}
        \end{definicao}
    \end{frame}

    \begin{frame}
        Podemos também escrever
        \begin{enumerate}[label=({\roman*})]
          \item $\proj_{\vec{u}}\vec{v} \varparallel\vec{u}$
            \item $(\vec{v} - \proj_{\vec{u}}\vec{v}) \perp \vec{u}$.
        \end{enumerate}

        No caso acima temos
        \[
          \cos\theta = \dfrac{\norm{\vec{OP}}}{\norm{\vec{v}}}
        \]
        e então $\norm{\vec{OP}} = \norm{\vec{v}}\cos\theta$. Mas
        \[
          \cos\theta = \dfrac{\inner{u}{v}}{\norm{\vec{u}}\norm{\vec{v}}}
        \]
        e assim podemos escrever
        \[
          \norm{\vec{OP}} = \dfrac{\inner{u}{v}}{\norm{\vec{u}}}.
        \]
    \end{frame}

    \begin{frame}
        Além disso, como $\vec{OP}$ é paralelo à $\vec{u}$ e de mesmo sentido, devemos ter $\vec{OP} = \lambda\vec{u}$ com $\lambda > 0$. Daí
        $\lambda = \norm{\vec{OP}}/\norm{\vec{u}}$ e então
        \[
          \vec{OP} = \dfrac{\norm{OP}}{\norm{\vec{u}}^2}\vec{u}.
        \]
    \end{frame}

    \begin{frame}
        Esta construção só é válida no caso em que o ângulo $\theta$ entre $\vec{u}$ e $\vec{v}$ satisfaz $0 < \theta < \pi/2$. Nos casos
        \begin{figure}
          \centering
          \caption{Vetor projeção}
          \includegraphics[scale=0.8]{vetores-projecao-plano-cartesiano.pdf}
        \end{figure}
        não podemos usar a construção envolvendo o ângulo $\theta$ entre os vetores.
    \end{frame}

    \begin{frame}
        \begin{proposicao}
          Seja $\vec{u}$ um vetor não nulo. Qualquer que seja o vetor $\vec{v}$, existe e é única a projeção ortogonal de $\vec{v}$ sobre $\vec{u}$. Sua expressão em termos de $\vec{u}$ e $\vec{v}$ é dada por
          \[
            \proj_{\vec{u}}\vec{v} = \dfrac{\inner{v}{u}}{\norm{\vec{u}}^2}\vec{u}
          \]
          e sua norma é dada por
          \[
            \norm{\proj_{\vec{u}}\vec{v}} = \dfrac{\mid\inner{v}{u}\mid}{\norm{\vec{u}}}.
          \]
        \end{proposicao}
    \end{frame}

    \begin{frame}
        \begin{observacao}
          O escalar $\lambda$ encontrado na proposição anterior pode ser negativo ou zero. Ele será negativo quando $\pi/2 \le \theta \le \pi$ e será zero quando $\vec{u}\perp\vec{v}$.
          \begin{figure}[!h]
            \centering
            \caption{Possibilidades para $\lambda$}
            \includegraphics{vetores-projecao-escalar-plano-cartesiano.pdf}
          \end{figure}
        \end{observacao}
    \end{frame}

    \begin{frame}
        \begin{exemplos}\label{exemplosprojecao}
          \begin{enumerate}[label=({\roman*})]
            \item Encontre a projeção ortogonal do vetor $\vec{v} = (4, -1)$ sobre o vetor $\vec{u} = (2, 1)$.
             \item Considere um triângulo $ABC$ onde $A(3,3)$, $B(0,1)$ e $C(1,6)$.
             \begin{enumerate}[label=({\roman*})]\label{exemploprojecao2}
               \item Verifique que $ABC$ é um triângulo retângulo em $A$.
               \item Calcule a projeção do cateto $BA$ sobre a hipotenusa $BC$.
               \item Determine o pé da altura do triângulo relativo ao vértice $A$.
             \end{enumerate}
          \end{enumerate}
        \end{exemplos}
    \end{frame}
\end{document}
